\documentclass[8pt,a4paper,twocolumn]{article}

\usepackage[spanish]{babel}
\usepackage[utf8]{inputenc}
\usepackage[T1]{fontenc}
\usepackage[left=0.75cm,right=0.75cm,top=0.55cm,bottom=0.7cm,columnsep=0.45cm]{geometry}
\usepackage{lmodern}
\usepackage{booktabs}
\usepackage{longtable}
\usepackage{array}
\usepackage{enumitem}
\usepackage{xcolor}
\usepackage{hyperref}
\usepackage{titlesec}
\usepackage[most]{tcolorbox}

\setlength{\parindent}{0pt}
\setlength{\parskip}{0.12em}
\setlength{\columnseprule}{0pt}
\pagestyle{plain}
\setlist[itemize]{leftmargin=1.0cm,itemsep=0.1em,topsep=0.15em}
\setlist[enumerate]{leftmargin=1.0cm,itemsep=0.1em,topsep=0.15em}
\titlespacing*{\section}{0pt}{0.35em}{0.2em}

\newcommand{\resumentitulo}{Resumen de estudio - Semana 14}

\begin{document}

\vspace*{-1.0em}
\begin{tcolorbox}[colback=black!12,colframe=black!35,boxrule=0.3pt,arc=1pt,left=4pt,right=4pt,top=3pt,bottom=3pt]
{\bfseries\normalsize \resumentitulo\par}
{\scriptsize CIB12 - Aprendizaje de Maquina y Mineria de Datos\par}
\end{tcolorbox}
\vspace{0.15em}

\section*{Tema}
Fundamentos de los algoritmos de aprendizaje profundo para modelos grandes

\section*{Indicaciones}
Reemplaza este contenido con el material generado por la IA o pega aqui el resumen final en LaTeX.

\section*{Resumen de la semana}
Escribe aqui una sintesis general del tema.

\section*{Conceptos clave}
\begin{itemize}
    \item Concepto 1
    \item Concepto 2
    \item Concepto 3
\end{itemize}

\section*{Desarrollo del tema}
Escribe aqui la explicacion principal.

\section*{Preguntas probables}
\begin{enumerate}
    \item Pregunta probable 1.
    \item Pregunta probable 2.
    \item Pregunta probable 3.
\end{enumerate}

\end{document}
